%%%%%%%%%%%%%%%%%%%%%%%%%%%%%%%%%%%%%%%%%%%%%%%%%%%%%%%%%%%%%%%%%%%%%%%%%%%%%%%%%%%%%%
%%%%%%%%%%%%%%%%%%%%%%%%%%%%%%%%%% Settings %%%%%%%%%%%%%%%%%%%%%%%%%%%%%%%%%%%%%%%%%%
%%%%%%%%%%%%%%%%%%%%%%%%%%%%%%%%%%%%%%%%%%%%%%%%%%%%%%%%%%%%%%%%%%%%%%%%%%%%%%%%%%%%%%

\documentclass{beamer}
\usepackage[utf8]{inputenc}
\usepackage[english]{babel}
\usepackage[T1]{fontenc}
\usepackage{appendixnumberbeamer}
\usepackage{siunitx}

\usepackage{minted}
\setminted{linenos, numberblanklines=false, mathescape, texcomments, autogobble, breakanywhere, breakautoindent, breaklines, frame=lines, framesep=5\fboxsep}
\setminted[python]{python3}

\usefonttheme{professionalfonts}
\usepackage{mathspec}
\setsansfont[BoldFont={Fira Sans},
Numbers={OldStyle}]{Fira Sans Light}
\setmathsfont(Digits)[Numbers={Lining, Proportional}]{Fira Sans Light}

\def\mytitle{Production Line Performance}
\def\mysubtitle{Project work in Languages and Algorithms for Artificial Intelligence class (Module 2)}
\def\myshortfullname{A. Falai}
\def\myfullname{Alessio Falai}
\def\myemail{alessio.falai@studio.unibo.it}
\def\myshortinstitute{UNIBO}
\def\myinstitute{Alma Mater Studiorum - University of Bologna}
\title{\mytitle}
\subtitle{\mysubtitle}
\author[\myshortfullname]{\myfullname \newline \texttt{\myemail}}
\date{\today}
\institute[\myshortinstitute]{\myinstitute}
\titlegraphic{\hfill\includegraphics[height=1.5cm]{img/unibo-logo.eps}}

\usetheme[progressbar=frametitle, numbering=none, background=light, titleformat=smallcaps, block=fill]{metropolis}
\useoutertheme{metropolis}
\useinnertheme{metropolis}
\usefonttheme{metropolis}
\usecolortheme{metropolis}
\setbeamercovered{dynamic}

% Show 'page' before page numbers in toc
\addtocontents{toc}{~\hfill page\par}

% Show page numbers and dots in toc
\makeatletter
\long\def\beamer@section[#1]#2{%
  \beamer@savemode%
  \mode<all>%
  \ifbeamer@inlecture
    \refstepcounter{section}%
    \beamer@ifempty{#2}%
    {\long\def\secname{#1}\long\def\lastsection{#1}}%
    {\global\advance\beamer@tocsectionnumber by 1\relax%
      \long\def\secname{#2}%
      \long\def\lastsection{#1}%
      \addtocontents{toc}{\protect\beamer@sectionintoc{\the\c@section}{#2\dotfill\the\c@page}{\the\c@page}{\the\c@part}%
        {\the\beamer@tocsectionnumber}}}%
    {\let\\=\relax\xdef\sectionlink{{Navigation\the\c@page}{\noexpand\secname}}}%
    \beamer@tempcount=\c@page\advance\beamer@tempcount by -1%
    \beamer@ifempty{#1}{}{%
      \addtocontents{nav}{\protect\headcommand{\protect\sectionentry{\the\c@section}{#1}{\the\c@page}{\secname}{\the\c@part}}}%
      \addtocontents{nav}{\protect\headcommand{\protect\beamer@sectionpages{\the\beamer@sectionstartpage}{\the\beamer@tempcount}}}%
      \addtocontents{nav}{\protect\headcommand{\protect\beamer@subsectionpages{\the\beamer@subsectionstartpage}{\the\beamer@tempcount}}}%
    }%
    \beamer@sectionstartpage=\c@page%
    \beamer@subsectionstartpage=\c@page%
    \def\insertsection{\expandafter\hyperlink\sectionlink}%
    \def\insertsubsection{}%
    \def\insertsubsubsection{}%
    \def\insertsectionhead{\hyperlink{Navigation\the\c@page}{#1}}%
    \def\insertsubsectionhead{}%
    \def\insertsubsubsectionhead{}%
    \def\lastsubsection{}%
    \Hy@writebookmark{\the\c@section}{\secname}{Outline\the\c@part.\the\c@section}{2}{toc}%
    \hyper@anchorstart{Outline\the\c@part.\the\c@section}\hyper@anchorend%
    \beamer@ifempty{#2}{\beamer@atbeginsections}{\beamer@atbeginsection}%
  \fi%
  \beamer@resumemode}%
\makeatother

% Change shape size for itemize
\setbeamertemplate{itemize item}[circle]
\setbeamertemplate{itemize subitem}[square]

% Conversion to roman numbers
\makeatletter
\newcommand*{\rom}[1]{\expandafter\@slowromancap\romannumeral #1@}
\makeatother


%%%%%%%%%%%%%%%%%%%%%%%%%%%%%%%%%%%%%%%%%%%%%%%%%%%%%%%%%%%%%%%%%%%%%%%%%%%%%%%%%%%%%%
%%%%%%%%%%%%%%%%%%%%%%%%%%%%%%%%%% Document %%%%%%%%%%%%%%%%%%%%%%%%%%%%%%%%%%%%%%%%%%
%%%%%%%%%%%%%%%%%%%%%%%%%%%%%%%%%%%%%%%%%%%%%%%%%%%%%%%%%%%%%%%%%%%%%%%%%%%%%%%%%%%%%% 

\begin{document}

%%%%%%%%%%%%%%%%%%%%%%%%%%%%%%%%%%%%%%%%%%%%%%%%%%%%%%%%%%%%%%%%%%%%%%%%%%%%%%%%%%%%%%
%%%%%%%%%%%%%%%%%%%%%%%%%%%%%%%%% Title & TOC %%%%%%%%%%%%%%%%%%%%%%%%%%%%%%%%%%%%%%%%
%%%%%%%%%%%%%%%%%%%%%%%%%%%%%%%%%%%%%%%%%%%%%%%%%%%%%%%%%%%%%%%%%%%%%%%%%%%%%%%%%%%%%% 

\renewcommand*{\insertpagenumber}{%
	\Roman{framenumber}%
}%
\makeatother
\setbeamertemplate{footline}
{
  \leavevmode%
  \hbox{%
  \begin{beamercolorbox}[wd=.3\paperwidth,ht=2.25ex,dp=1ex,center]{author in head/foot}%
    \usebeamerfont{author in head/foot}\insertshortauthor ~ (\myshortinstitute)
  \end{beamercolorbox}%
  \begin{beamercolorbox}[wd=.6\paperwidth,ht=2.25ex,dp=1ex,center]{block title}%
    \usebeamerfont{title in head/foot}\insertshorttitle
  \end{beamercolorbox}%
  \begin{beamercolorbox}[wd=.1\paperwidth,ht=2.25ex,dp=1ex,center]{block body}%
    \insertpagenumber{} \hspace*{1ex}
  \end{beamercolorbox}}%
  \vskip0pt%
}
\makeatletter
\setbeamertemplate{navigation symbols}{}

\begin{frame}
	\maketitle
\end{frame}

\begin{frame}
	\frametitle{Table of contents}
	\setbeamertemplate{section in toc}[ball]
	\tableofcontents[hideallsubsections]
\end{frame}

%%%%%%%%%%%%%%%%%%%%%%%%%%%%%%%%%%%%%%%%%%%%%%%%%%%%%%%%%%%%%%%%%%%%%%%%%%%%%%%%%%%%%%
%%%%%%%%%%%%%%%%%%%%%%%%%%%%%%%%% Main content %%%%%%%%%%%%%%%%%%%%%%%%%%%%%%%%%%%%%%%
%%%%%%%%%%%%%%%%%%%%%%%%%%%%%%%%%%%%%%%%%%%%%%%%%%%%%%%%%%%%%%%%%%%%%%%%%%%%%%%%%%%%%% 

\setcounter{framenumber}{0}
\makeatother
\setbeamertemplate{footline}
{
  \leavevmode%
  \hbox{%
  \begin{beamercolorbox}[wd=.3\paperwidth,ht=2.25ex,dp=1ex,center]{author in head/foot}%
    \usebeamerfont{author in head/foot}\insertshortauthor ~ (\myshortinstitute)
  \end{beamercolorbox}%
  \begin{beamercolorbox}[wd=.6\paperwidth,ht=2.25ex,dp=1ex,center]{block title}%
    \usebeamerfont{title in head/foot}\insertshorttitle
  \end{beamercolorbox}%
  \begin{beamercolorbox}[wd=.1\paperwidth,ht=2.25ex,dp=1ex,center]{block body}%
    \insertframenumber{} / \inserttotalframenumber\hspace*{1ex}
  \end{beamercolorbox}}%
  \vskip0pt%
}
\makeatletter
\setbeamertemplate{navigation symbols}{}

% Datasets section
\section{Datasets}
\begin{frame}
  \frametitle{Test datasets}
  \begin{itemize}[<+->]
    \item \textbf{Arrest}: Contains statistics, in arrests per 100.000 residents, for assault, murder, and rape in each of the 50 US states in 1973. 
          It also gives the percent of the population living in urban areas.
    \item \textbf{Adult}: Aims at separating people whose income is greater than 50 thousands dollars per year from the rest.
  \end{itemize}	
\end{frame}

\begin{frame}
  \frametitle{Actual dataset}
  \begin{itemize}[<+->]
    \item \textbf{Bosch}: Aims at predicting internal failures using thousands of measurements and tests made for each component along assembly lines.
  \end{itemize}
\end{frame}

% Preprocessing section
\section{Preprocessing}
\begin{frame}
  \frametitle{Data exploration}
  \begin{itemize}[<+->]
    \item The Bosch dataset includes 3 subsets: numerical, categorical and time data.
    \item As stated in \cite{predict-failures}, the categorical data is extremely sparse, and thus not exploited in subsequent stages.
    \item Our first analysis will be focused on just numerical data:
    \begin{itemize} 
      \item $968$ anonymized features
      \item<.-> \num{1183747} labeled examples
      \item<.-> $0.58\%$ of failed products
      \item<.-> $78.5\%$ of missing values
    \end{itemize}
  \end{itemize}
\end{frame}

\begin{frame}
  \frametitle{Two-stage approach}
  \begin{itemize}[<+->]
    \item \textbf{Stage \rom{1}}: This step clusters data with similar processes together into process groups.
    \item \textbf{Stage \rom{2}}: This step uses supervised learning to predict the failed products. Each cluster is treated as an independent dataset and has its own classifier.
  \end{itemize}
\end{frame}

\begin{frame}
  \frametitle{Data preprocessing}
  \begin{itemize}[<+->]
    \item \textbf{Common}: Columns containing null values and constant values in each row are dropped.
    \item \textbf{Clustering}: Values get binarized (0 meaning null value and 1 meaning not null value) and PCA is applied to binarized features.
    \item \textbf{Classification}: A feature imputation method (mean value over the column), followed by feature standardization (zero mean, unit variance) and PCA, is applied over non-binary values in each cluster.
    \item \textbf{Prediction}: A new example follows the same preprocessing scheme as the whole dataset.
  \end{itemize}
\end{frame}

% Clustering section
\section{Clustering}
\begin{frame}
  \frametitle{k-means}
    The chosen clustering algorithm is k-means, since density-based methods (like DBSCAN) are still not available in Spark. So, the following issues needed to be addressed:

    \begin{itemize}
      \item \textbf{Problem}: Automatically select the right amount of clusters:
      \begin{itemize}
        \item Silhouette and elbow methods are not ideal, since they require human analysis.
        \item Spark 3.0.0 dropped support for inertia computation, maintaining only evaluation by silhouette scores.
      \end{itemize} 
      \item \textbf{Solution}: Ad-hoc implementation of the Gap statistic method, described in \cite{gap}.
    \end{itemize}
\end{frame}

\begin{frame}
  \frametitle{Gap statistic}
    \begin{itemize}
      \item For $k = 1, .., K$, where $K$ is the maximum number of clusters, perform k-means and compute the resulting inertia $I_{k}$.
      \item Generate $B$ reference datasets by sampling from a uniform distribution over each feature, where the support is directly identified by features ranges. Then, for $k = 1, .., K$ and $b = 1, .., B$, perform k-means and compute the resulting inertia $I_{kb}$. Finally, estimate $E^{*}[log(I_{kb})]$ as $\frac{1}{B} \sum_{b}log(I_{kb})$.
      \item Compute the Gap score as $Gap(k) = E^{*}[log(I_{kb})] - log(I_{k})$.
      \item Compute the standard deviation $sd_{k}$ of $log(I_{kb})$ and define $s_{k} = sd_{k} * \sqrt{1 + \frac{1}{B}}$.
      \item Select the minimum $k$ s.t. $Gap(k) - Gap(k + 1) + s_{k + 1} \geq 0$.
    \end{itemize}
\end{frame}

% Classification section
\section{Classification}
\begin{frame}
  
\end{frame}

% Cloud section
\section{Cloud}
\begin{frame}

\end{frame}

%%%%%%%%%%%%%%%%%%%%%%%%%%%%%%%%%%%%%%%%%%%%%%%%%%%%%%%%%%%%%%%%%%%%%%%%%%%%%%%%%%%%%%
%%%%%%%%%%%%%%%%%%%%%%%%%%%%%%%%%%% Appendix %%%%%%%%%%%%%%%%%%%%%%%%%%%%%%%%%%%%%%%%%
%%%%%%%%%%%%%%%%%%%%%%%%%%%%%%%%%%%%%%%%%%%%%%%%%%%%%%%%%%%%%%%%%%%%%%%%%%%%%%%%%%%%%% 

\appendix

\begin{frame}[standout]
	Thank you for your attention
\end{frame}

\begin{frame}
	\frametitle{References}
	\nocite{*}
    \bibliography{bibliography}
    \bibliographystyle{plain}
\end{frame}

\begin{frame}
	\frametitle{Backup frame}
	\usebeamercolor[fg]{normal text}
	This is a backup frame, useful to include additional material for questions from the audience.
\end{frame}
\end{document}
